\section{From Current Decisions to Policy Accuracy}\label{sec:r21composition}
A continuation certificate is useful for a dynamic economy only if the
certified object is the future of the policy under assessment. This section
states the conditions under which the local decision accounts compose.
They apply to additive and order-preserving recursive evaluation, while
keeping the utility domains and admissible deviations of the economic model.

At dates $t=0,\ldots,T-1$, let $\mathcal D_t$ be an invariant state domain
and $A_t(s)$ the full admissible current action set. Write
\[
 (T_t v)(s)=\sup_{a\in A_t(s)}G_t(s,a,v),\qquad
 (T_t^\pi v)(s)=G_t(s,\pi_t(s),v).
\]
Assume the dynamic programming principle on these domains, measurable
admissible policies, and finite values. Each $G_t$ is order preserving in its
continuation argument and is sup-norm Lipschitz, uniformly in $(s,a)$, with
constant $\beta_t\geq0$ on a common invariant value domain. Let
$V_T^*=V_T^\pi$ and define $V_t^*=T_tV_{t+1}^*$ and
$V_t^\pi=T_t^\pi V_{t+1}^\pi$. The relevant local advantage is
\begin{equation}
 g_t^\pi(s)=T_tV_{t+1}^\pi(s)-T_t^\pi V_{t+1}^\pi(s).
 \label{eq:r21advantage}
\end{equation}
It evaluates all current deviations against the \emph{same implemented
future policy}; it is neither a training residual nor an advantage computed
against an unrelated anchor continuation.

\begin{theorem}[Composition of economic certificates]\label{thm:r21composition}
Under the preceding assumptions, suppose that on a common event
\begin{equation}
 0\leq g_t^\pi(s)\leq\varepsilon_t
 \quad\text{for every }s\in\mathcal D_t,
 \qquad t=0,\ldots,T-1.
 \label{eq:r21uniform}
\end{equation}
Then, on that event, at every initial state,
\begin{equation}
 0\leq V_0^*-V_0^\pi\leq
 \sum_{t=0}^{T-1}w_t\varepsilon_t,
 \qquad w_t=\prod_{j=0}^{t-1}\beta_j.
 \label{eq:r21composition}
\end{equation}
For additive Markov evaluation
$G_t(s,a,v)=r_t(s,a)+\beta_tP_t^av(s)$, a state-dependent refinement is
$V_t^*-V_t^\pi\leq B_t$, where
\begin{equation}
 B_T=0,\qquad
 B_t(s)=u_t(s)+\beta_t\sup_{a\in A_t(s)}P_t^aB_{t+1}(s)
 \label{eq:r21envelope}
\end{equation}
whenever $u_t(s)\geq g_t^\pi(s)$ and the displayed quantities are finite.
\end{theorem}

Section~\ref{app:r21proofs} proves the result. Centered continuation-error
and transport bounds from Proposition~\ref{prop:r19decision} can supply
\eqref{eq:r21uniform} when established uniformly on the required state and
action domains. Alternatively, an interval certificate for the true
one-step objective can supply it directly. In either case, the continuation
must be $V_{t+1}^\pi$, or its discrepancy must be charged. The verification
may be deterministic. With random certificates, all domain and date bounds
must hold on one event, for example by a prespecified summable conditional
error allocation. Selecting a stopping time does not remove this requirement.

This theorem identifies a positive route from reusable policy evaluation to
full policy accuracy. It also identifies precisely what the computational
evidence below supplies. A mean certificate on a finite state catalogue
establishes the declared current-decision experiment; it does not establish
\eqref{eq:r21uniform} outside that catalogue. Sampling only states visited by
$\pi$ is insufficient because a profitable deviation may reach other states.
The state envelope in \eqref{eq:r21envelope} deliberately takes a supremum
over admissible transitions, not only the transitions of the fitted policy.

\paragraph{Relation to the continuous economy and strategic applications.}
Let $\Pi_h$ be a specified implementable policy class in the original economy.
If its class-approximation gap is at most $\eta_{\rm class}$, and
$|J(\pi)-J_h(\pi)|\leq\eta_{\rm value}$ uniformly for $\pi\in\Pi_h$, then a
policy satisfying \eqref{eq:r21composition} in the approximating model has
\begin{equation}
 \sup_{\pi}J(\pi)-J(\widehat\pi)
 \leq\eta_{\rm class}+2\eta_{\rm value}
       +\sum_t w_t\varepsilon_t+\eta_{\rm impl},
 \label{eq:r21bridge}
\end{equation}
where $\eta_{\rm impl}$ bounds any additional payoff loss in implementation.
The class and value-transfer allowances are separate mathematical inputs,
not numbers inferred from the finite experiments. For recursive utility,
order and Lipschitz bounds are imposed on the actual common utility domain;
actionwise centering is applied to the evaluated action payoff, not through
an arbitrary nonlinear aggregator. For a dynamic game, apply the same
account to each player's full deviation problem against the same opponents'
profile. Bounds $\epsilon_i$ for all players imply an $\epsilon_i$-best-response
profile under the specified equilibrium concept. A single local deviation
or a single player's calculation does not supply those premises. The original
information, boundary, transversality, and deviation requirements remain in
the Economic Applications companion.

\subsection{Allocating computation across dates}
The policy target suggests how to distribute a verified approximation budget.
Suppose the local account at date $t$ is $b_t+e_t$, where $b_t$ contains
fixed optimization, representation, and transport allowances. Suppose further
that a declared procedure attains approximation allowance $e_t>0$ using
$n_t\geq A_t/e_t^2$ fitting observations at charged marginal cost $c_t>0$.
Here $A_t>0$ includes its confidence allocation and design constants. The
full-rank and coverage conditions underlying such a rate must hold at the
selected budgets; Corollary~\ref{cor:r19amortization} is one finite-query
example, not a license to assume uniform coverage of an infinite state space.

\begin{proposition}[Allocation for a policy target]\label{prop:r21allocation}
Assume $w_t>0$ and let
$E=\varepsilon-\sum_t w_tb_t>0$, $a_t=c_tA_t$, and
$H=\sum_ta_t^{1/3}w_t^{2/3}$. Among positive allowances satisfying
$\sum_tw_te_t\leq E$, the minimum continuous fitting bill is
\begin{equation}
 \min\sum_t\frac{a_t}{e_t^2}=\frac{H^3}{E^2},\qquad
 e_t=\frac{E}{H}\left(\frac{a_t}{w_t}\right)^{1/3}.
 \label{eq:r21allocation}
\end{equation}
Rounding each observation count up adds at most $\sum_tc_t$ to this bill.
All fixed construction, reusable-cache, query, and verification costs must
be added. The statement minimizes the displayed sufficient-cost account,
not the work of all possible numerical methods.
\end{proposition}

The propagation and allocation arguments use standard dynamic-programming
stability and convex optimization. Their role here is to connect the paper's
centered, transported, and true-objective certificates to a specified dynamic
economic target. They do not replace approximation or coverage assumptions.
Finite-time fitted-value analysis, including the distinction between sampling
norms and policy loss, provides an important antecedent
\citep{MunosSzepesvari2008}. The empirical contribution below remains a
measured comparison of complete procedures, rather than a claim of a new
principle of optimality.
